%% SCRAP: architecture/03-architecture/hal/starkernel-integration
%% SOURCE: docs/working/architecture/03-architecture/hal/starkernel-integration.md
%% STATUS: WORKING
%% FITS: dev-guide/ch-platform
%% EDITORIAL: lifted — prose rewritten to press voice

\section{StarKernel HAL Integration}

This section gives the kernel-specific detail for the HAL platform layer that
becomes StarKernel: the UEFI boot sequence, the freestanding C environment,
hardware bring-up, and the path to a working \texttt{ok} prompt.

\subsection{Boot Architecture}

The kernel boots in four handoffs. UEFI firmware initializes hardware, exposes
boot services --- the memory map, the ACPI tables, and the graphics output
protocol --- and loads the boot image. The StarKernel boot loader collects that
information into a \texttt{BootInfo} structure, sets up initial page tables,
exits boot services, and jumps to the kernel proper. The kernel HAL stage
initializes the CPU structures, the memory subsystem, the time subsystem, and the
console. Finally the StarForth VM is created, its physics subsystems are
initialized, the REPL starts, and the \texttt{ok} prompt appears.

\subsection{UEFI Loader}

The loader is the UEFI entry point. It collects the memory map (over-allocating
to leave room for the descriptor that \texttt{ExitBootServices} will add), locates
the ACPI 2.0 RSDP among the configuration tables, queries the graphics output
protocol for the framebuffer base and geometry, exits boot services, and transfers
control to the kernel. It never returns.

\begin{lstlisting}[language=C]
typedef struct {
    uint64_t memory_map_addr;
    uint64_t memory_map_size;
    uint64_t memory_map_descriptor_size;
    uint64_t acpi_rsdp_addr;
    uint64_t framebuffer_addr;
    uint32_t framebuffer_width;
    uint32_t framebuffer_height;
    uint32_t framebuffer_pitch;
} BootInfo;
\end{lstlisting}

\subsection{Kernel Entry Point}

The kernel entry point parses \texttt{BootInfo} and brings the machine up in
order: it initializes the serial console first so that every later step can emit
debug output, then the descriptor tables, then the physical and virtual memory
managers, then the kernel heap, then the framebuffer, and finally the HAL
subsystems. With the HAL ready it calls the standard \texttt{main()} to start the
VM, and panics if \texttt{main()} ever returns.

\subsection{HAL Implementation Notes}

The kernel HAL rests on a handful of hardware facilities. Time calibrates the
TSC against the HPET, checks for an invariant TSC, records a boot timestamp, and
drives periodic interrupts from the local APIC timer.

\begin{lstlisting}[language=C]
uint64_t hal_time_now_ns(void) {
    uint64_t tsc = rdtsc() - boot_tsc;
    return (tsc * 1000000000ULL) / tsc_hz;
}
\end{lstlisting}

Interrupts populate a 256-entry IDT, mask the legacy PIC in favor of the local
APIC and IOAPIC, and track nesting depth so that \texttt{hal\_in\_interrupt\_context()}
is exact. Memory layers a bitmap physical allocator over the parsed UEFI memory
map, a four-level page-table virtual manager, and a \texttt{kmalloc} heap that
zero-initializes to honor the contract. The console drives a 16550 UART at
115200 baud in polling mode and, where present, a linear GOP framebuffer with
software text rendering. CPU support reports the local APIC identity and exposes
the relax and halt primitives.

\subsection{Build System}

The kernel is built with a freestanding toolchain: a C compiler invoked with
\texttt{-ffreestanding} and floating point disabled, a linker driven by a custom
script, and \texttt{objcopy} to emit a PE32+ executable for UEFI. The linker
script loads the image at the one-megabyte mark and lays out text, read-only
data, data, and BSS while discarding the exception-handling frame sections.

\subsection{Testing on QEMU}

The kernel is exercised under QEMU with OVMF as the UEFI firmware. The image is
copied into an EFI system partition as \texttt{BOOTX64.EFI} and booted with a
serial console routed to standard output. A healthy boot prints its progress
through the memory, heap, and HAL initialization stages and ends at the VM banner
and the \texttt{ok} prompt.

\begin{lstlisting}[language=bash]
qemu-system-x86_64 \
    -bios /usr/share/ovmf/OVMF.fd \
    -drive file=fat:rw:esp/,format=raw \
    -serial stdio -m 512M -enable-kvm
\end{lstlisting}

\subsection{Debugging}

Three techniques carry most of the load. Early serial output writes directly to
the UART before any higher-level console exists. GDB attaches to QEMU's remote
stub for source-level debugging from the entry point onward. The panic handler
disables interrupts, prints its message to the console, and halts every CPU.

\subsection{Roadmap to the \texttt{ok} Prompt}

The bring-up advances through seven milestones: boot with serial output and no
triple fault; working memory with a functioning physical manager, virtual
mapping, and heap; a fully initialized HAL; a VM that creates and allocates its
dictionary without crashing; a REPL that prints \texttt{ok}, echoes input, and
executes simple words; operational physics subsystems with a firing heartbeat;
and finally the full test suite passing on the kernel with zero algorithmic
variance on bare metal.

%% TODO(bob): reconcile milestone numbering here (M1--M7) with the LithosAnanke
%%   M0--M7 milestones in the kernel overview; confirm current completion state.

Beyond the prompt lies StarshipOS: storage drivers and a filesystem, networking,
a Forth-task process model, a unified device model, and a capability- and
ACL-based security model --- all built, as before, on the HAL.
